% !TEX root = ../main.tex
% ============================================================
% Discussion chapter
% Revised after the 2026-05-21 discussion with Stephen Johnson and
% the 2026-05-18 professor overview/results notes.
%
% Style constraints:
% - No personal pronouns.
% - No paragraph-style headings.
% - No em-dashes.
% - Booktabs only for tables.
% - Quantitative claims are taken from the current Results chapter,
%   local thesis notes, or meeting notes.
% ============================================================

\chapter{Discussion}
\label{chapter:discussion}

The preceding results support a deliberately narrower conclusion than the
initial motivation of the project. Learned frequency transfer can predict
high-frequency Helmholtz fields, and full-grid FD/PML training can improve the
finite-budget true-residual trajectory of CSL-FGMRES. The learned model is not
yet a competitive replacement for classical Helmholtz solvers, and the present
evidence does not justify calling the current two-dimensional model a learned
preconditioner. The scientific contribution is instead the identification of a
solver-facing mechanism: field error becomes relevant to Krylov convergence
only through the residual that the discrete operator induces.

The central object is therefore not the neural network itself but the
operator-amplified residual
\begin{equation}
    r_0 = b - A_H x_0 = A_H e_0,
    \qquad
    e_0 = u_H - x_0 .
    \label{eq:discussion-error-residual}
\end{equation}
In the one-dimensional Dirichlet eigenbasis this becomes
\begin{equation}
    c_k(r_0) = \lambda_k c_k(e_0),
    \label{eq:discussion-modal-amplification}
\end{equation}
as derived in Section~\ref{sec:ch5-residual-amplification} and measured in
Section~\ref{sec:results-1d}. Equation~\eqref{eq:discussion-modal-amplification}
explains why a field-accurate warm start can be a poor algebraic initial
state. It also explains why residual-aware learning, residual gating, and
operator-compatible PML treatment are not optional refinements but direct
consequences of the solver interface.

The discussion is organised around seven questions. Section~\ref{sec:discussion-synthesis}
synthesises the one-dimensional Dirichlet, one-dimensional PML, and
two-dimensional FD/PML evidence. Section~\ref{sec:discussion-what-works}
states what the current method achieves. Section~\ref{sec:discussion-what-fails}
states what it does not achieve. Section~\ref{sec:discussion-classical}
positions the work against classical Helmholtz alternatives, including the
Green's-function and low-rank baselines that are especially strong in the
linear homogeneous case. Section~\ref{sec:discussion-amortisation} discusses
training cost and the application regimes in which the cost could be
amortised. Section~\ref{sec:discussion-limitations} records threats to
validity. Section~\ref{sec:discussion-future} gives the research programme
that follows most directly from the results.

% ============================================================
\section{Synthesis of the evidence}
\label{sec:discussion-synthesis}
% ============================================================

The empirical record across the three solver settings is best read as one
phenomenon observed under progressively fewer simplifying assumptions. The
Dirichlet experiment supplies the exact modal identity
\eqref{eq:discussion-modal-amplification}. The one-dimensional PML experiment
removes the self-adjoint eigenbasis and tests whether the same residual
diagnostics survive when the absorbing layer is part of the operator. The
two-dimensional FD/PML experiment is the deployment setting, with a
complex-valued, non-normal operator on a $512\times512$ grid.

Three results are directly predicted by the theory.

\textit{Field accuracy is not residual accuracy.} In the 1D Dirichlet
$16\to32$ problem, the raw U-Net has field error $0.071$ but true initial
residual $12.81\|b\|_2$, while the cold start has true initial residual
$\|b\|_2$ by construction. This is the clearest numerical instance of
\eqref{eq:discussion-modal-amplification}. The field error is small in the
solution norm, but its modal content is not residual-safe.

\textit{CSL changes the residual view.} The same raw 1D Dirichlet U-Net start
has CSL-preconditioned residual $0.173$, below the cold-start value. This does
not contradict the large true residual. It shows that the CSL preconditioner
filters the residual in a way that is more forgiving to some of the modes that
the original Helmholtz operator amplifies. This distinction explains why both
true and preconditioned residuals must be reported.

\textit{PML compatibility is operator compatibility.} The absorbing layer is
part of the discrete operator used to evaluate $b-Ax$. The one-dimensional PML
experiment shows the extreme form of this fact: zeroing the PML output of the
older green-function model increases the true initial residual from $1.924$ to
$37.14$, while the flux-full model has true initial residual $0.539$. The newer
two-dimensional flux-full results show the same principle in a more useful
form: once the model is trained on the full FD/PML field, keeping the learned
PML output is better than zeroing it in the final true-residual metric for all
tested pairs.

Two results were not predicted quantitatively but are consistent with the same
framework.

The first is the modest size of the warm-start iteration gain. The residual
gate lowers the 1D Dirichlet true initial residual below the cold-start value,
but it does not significantly change the number of CSL-FGMRES iterations. A
warm start changes the initial residual vector; it does not change the
preconditioned operator or the polynomial approximation problem solved by
GMRES \cite{Saad2003}. This is why the method can improve the first part
of the convergence curve without changing the tail.

The second is the gap between exact coarse correction and learned correction.
In the 1D in-FGMRES diagnostic, CSL alone requires mean iteration count
$16.00$, while exact low solve plus $T_\uparrow$ with a residual gate reduces
the count to $12.65$. Current learned residual-correction variants remain at
approximately the CSL baseline. This is an important positive-negative result:
the solver-aligned mechanism exists, but current learned maps do not yet
reproduce it.

\begin{table}[htbp]
\centering
\caption{Solver-facing interpretation of the main evidence.}
\label{tab:discussion-evidence-map}
\begin{tabular}{lll}
\toprule
Setting & Main observation & Interpretation \\
\midrule
1D Dirichlet & Field error $0.071$, true residual $12.81$ & Modal residual amplification \\
1D Dirichlet gate & True residual lowered to $0.708$ & Useful and harmful modes coexist \\
1D exact correction & Iterations $16.00\to12.65$ & Solver-aligned coarse action exists \\
1D learned correction & Iterations stay near $16$ & Present learned target is insufficient \\
1D PML & Flux-full residual $0.539$ & PML must be learned as operator data \\
2D FD/PML & Flux-full raw gives best final true residual & Full-grid training helps solver trajectory \\
\bottomrule
\end{tabular}
\end{table}

Table~\ref{tab:discussion-evidence-map} gives the compact version. The thesis
does not show that a generic neural frequency-transfer model is a good
preconditioner. It shows that residual amplification is the reason field
transfer and solver acceleration diverge, and it identifies the solver-aligned
object that a successful learned correction must imitate.

% ============================================================
\section{What the current method achieves}
\label{sec:discussion-what-works}
% ============================================================

The present method is a learned warm start for CSL-FGMRES. Under that
definition it achieves four useful things.

\subsection{It learns nontrivial cross-frequency structure}
\label{subsec:discussion-cross-frequency}

The field-prediction results show that the learned models exploit information
in the low-frequency solution. The one-dimensional Dirichlet U-Net reaches a
small field loss in the clean proof setting. The two-dimensional flux-full
TransferUNet reaches full-grid validation RelL2 values of approximately
$0.240$, $0.321$, and $0.339$ for $16\to32$, $32\to64$, and $64\to128$,
respectively. These errors are far above the residual-safety threshold for a
general preconditioner, but they are far below the zero predictor and establish
that the cross-frequency map is learnable.

This point matters because it separates the negative solver results from a
trivial learning failure. The learned map is not useless. It captures
physically meaningful cross-frequency structure, but the captured structure is
not automatically aligned with the residual norm used by the solver.

\subsection{It improves finite-budget true residuals in 2D FD/PML}
\label{subsec:discussion-2d-success}

The latest two-dimensional beta-$0.3$ evaluation uses $10$ test right-hand
sides per frequency pair and $40$ CSL-FGMRES iterations. No method reaches the
tolerance within the budget, so the meaningful comparison is the residual
after the fixed budget. In that metric, the full-grid flux-trained warm start
has the lowest final true residual for all three pairs:
\[
    2.05\times10^{-3} \to 4.22\times10^{-4}
    \quad (16\to32),
\]
\[
    3.83\times10^{-1} \to 3.28\times10^{-1}
    \quad (32\to64),
\]
and
\[
    4.33\times10^{-1} \to 3.79\times10^{-1}
    \quad (64\to128).
\]
The gains are not dramatic at the harder pairs, but they are consistent in the
primary true-residual metric.

This is the strongest current two-dimensional claim. The method can be called
a useful warm start in the tested FD/PML setting, provided that ``useful'' is
understood as a finite-budget residual improvement rather than convergence
acceleration to tolerance.

\subsection{It shows that full-grid FD/PML training fixes part of the PML pathology}
\label{subsec:discussion-pml-fix}

The older interior-focused depth5 model had good physical-interior field
errors but harmful raw PML output. In the older N5/K20 snapshot, zeroing the
PML strip could turn a harmful learned field into a useful warm start for
$16\to32$. That conclusion no longer describes the best current model. The
newer flux-full model is trained against the full FD/PML solution, including
the absorbing layer, and keeping its PML output gives better final true
residuals than zeroing it.

The corrected interpretation is therefore not that the PML should be zeroed.
It is that the PML must be controlled. For an interior-only model, zeroing can
be a safety patch. For a full-grid FD/PML model, zeroing destroys part of the
operator-compatible output. This distinction is a useful methodological result
for learned Helmholtz solvers with absorbing boundaries.

\subsection{It produces an architecture-independent diagnostic}
\label{subsec:discussion-diagnostic}

The residual-amplification identity is independent of whether the field was
produced by a U-Net, a dilated CNN, a Fourier neural operator, or a classical
interpolant. Any approximate high-frequency field used as an initial guess
induces a residual through $A_H e_0$. The architecture affects the error
distribution, but the operator decides how that error is seen by the solver.

This is why the main contribution survives the modest iteration gains. The
diagnostic tells which learned outputs are solver-safe, why some accurate
fields are harmful, and what a residual-aware extension must target.

% ============================================================
\section{What the current method does not achieve}
\label{sec:discussion-what-fails}
% ============================================================

The same evidence imposes sharp limits on the claims.

\subsection{A warm start does not change the Krylov tail}
\label{subsec:discussion-warm-tail}

A warm start changes $r_0$ but not the spectrum or pseudospectrum of the
preconditioned operator. Once the initial residual has been processed, the
remaining convergence is controlled by the CSL-preconditioned FGMRES process.
This explains the repeated observation that early residual behaviour can
improve without a large iteration-count reduction. In practical terms, a warm
start can be expected to help the first few iterations, but the expensive tail
from moderate residual to tight tolerance remains a classical Krylov problem.

This point is especially important for the defence of the method. The present
warm-start result should not be advertised as a factor-of-ten solver
acceleration. It is an initial-state improvement and a diagnostic step toward
a residual-trained correction.

\subsection{The current learned correction does not imitate the exact coarse action}
\label{subsec:discussion-learned-correction-fails}

The exact-coarse diagnostic is the strongest evidence that a solver-aligned
frequency-transfer mechanism can reduce FGMRES work. It also reveals the
current gap. Learned residual-loss correction, scalar alpha rescaling, and
learned restriction through $T_\uparrow$ do not reproduce the useful modal
action of the exact correction. The learned restriction gate accepting almost
no modes is particularly informative: the learned map is not aligned with the
coarse correction that CSL-FGMRES can exploit.

The failure should not be interpreted as a generic failure of learning. It is
more specific. A solution-trained map, and even a generic residual-loss map,
does not necessarily learn the scale, phase, and modal selectivity of the
accepted exact coarse correction. The next training target must imitate that
solver-relevant action directly.

\subsection{The preconditioned residual is not uniformly solved}
\label{subsec:discussion-precond-residual}

The 2D flux-full warm start improves final true residuals, but the
preconditioned residual tells a more complicated story. For $64\to128$, the
raw full-grid learned start has initial preconditioned residual $18.41$ and
final preconditioned residual $28.51$, while still giving the best final true
residual. This means that the learned output is not uniformly compatible with
the coordinates seen by the CSL-preconditioned Krylov process.

This is the practical reason for retaining a cautious label. The method is a
warm start, not a learned preconditioner. A preconditioner claim would require
robust improvement in the residual norms relevant to the Krylov process and
would need to hold inside the iteration, not only at the initial state.

\subsection{High-frequency behaviour remains fragile}
\label{subsec:discussion-high-frequency}

The highest pair, $64\to128$, remains the stress test. The flux-full model has
the best final true residual after $40$ iterations, but the improvement over
cold start is modest and the preconditioned residual is unfavourable. This is
consistent with the threshold story: as frequency increases, the solver-safe
field-error budget becomes tighter, while fixed training objectives do not
automatically place the remaining error in benign modes.

The correct conclusion is not that high-frequency learned transfer is solved.
The correct conclusion is that the current high-frequency result confirms the
need for residual-aware training and conservative acceptance tests.

% ============================================================
\section{Comparison with classical alternatives}
\label{sec:discussion-classical}
% ============================================================

The strongest criticism of learned frequency transfer in the linear
homogeneous setting is not that the learned model fails to learn. It is that
classical methods already exploit linearity, superposition, and Green's
function structure extremely well. A fair discussion must therefore compare
against methods that use the same structure, not only against weak generic
preconditioners.

\subsection{Complex shifted Laplacian}
\label{subsec:discussion-csl}

CSL remains the backbone solver in this thesis. It changes the indefinite
Helmholtz problem into a shifted system by adding an imaginary damping term,
and the learned model only changes the initial guess supplied to the
CSL-preconditioned Krylov method \cite{Erlangga2004}. The present results are
therefore complementary to CSL rather than competitive with it. CSL shapes the
iteration; the learned model attempts to make the starting residual more
favourable.

This distinction should remain explicit. Any claim that the neural model
``beats CSL'' would be inaccurate, because CSL is still the preconditioner in
the successful experiments.

\subsection{Green's-function quadrature for the linear homogeneous problem}
\label{subsec:discussion-greens-quadrature}

For a fixed homogeneous linear Helmholtz operator, the solution is a
convolution with a Green's function. If the source is smooth or if the quantity
of interest is away from the source, a classical quadrature or spectral
expansion can approximate the solution from a modest number of precomputed
responses. A Chebyshev or Clenshaw-Curtis construction uses a small set of
quadrature-point solves, obtains coefficients by a cosine transform, and
forms new solutions by linear combination of the precomputed Green's-function
responses.

This baseline is much stronger than a cold-start iterative solve when its
assumptions hold. It preserves linearity exactly, needs far fewer than the
approximately $9600$ training solves used per pair in the current neural
datasets, and can converge rapidly for smooth source classes. In that regime,
the present learned model is not the most efficient way to approximate the
linear solution map.

The relevance of this comparison is conceptual as well as practical. Neural
networks can accidentally discard linear structure that the PDE gives for
free. The value of a learned component must therefore come from regimes where
linearity, smoothness, low dimensionality, or explicit Green's functions no
longer provide the dominant shortcut.

\subsection{Randomised low-rank Green's-function surrogates}
\label{subsec:discussion-low-rank}

Even when an explicit Green's function is unavailable, fixed linear wave
operators often have low-rank off-diagonal Green's-function blocks. Randomised
probing and truncated low-rank approximation can exploit this structure with a
number of solves proportional to the numerical rank rather than to the number
of possible right-hand sides. Hierarchical and sweeping methods rely on the
same broad fact: wave propagation has exploitable structure away from local
singular interactions \cite{Engquist2011}.

This is another reason the training-cost question is severe in the linear
case. A learned surrogate trained from thousands of full solves must be
amortised across a very large number of future solves, or it must handle a
regime where the low-rank linear surrogate is inadequate.

\subsection{Frequency-sweep and subspace-recycling methods}
\label{subsec:discussion-recycling}

The thesis motivation is cross-frequency reuse. Classical Krylov methods also
reuse information across nearby linear systems through subspace recycling,
deflation, and rational approximation. GCRO-DR and related recycled Krylov
methods use approximate invariant subspaces from earlier solves, while
rational or Padé approximants exploit analytic dependence on frequency for
families of shifted systems \cite{ParksEtAl2006}. These methods are natural
competitors for any linear frequency-transfer approach.

The present neural approach differs by learning a direct nonlinear map between
solution fields rather than recycling a linear subspace. That distinction is
not automatically an advantage in the linear setting. It becomes more
interesting when the parameter dimension grows, the source distribution is
non-smooth or widely scattered, or the governing equation becomes nonlinear.

\subsection{Where the learned approach may have a fair opening}
\label{subsec:discussion-fair-opening}

The classical baselines are strongest for fixed linear operators, smooth or
low-dimensional source families, far-field quantities of interest, and tight
accuracy requirements. The learned approach has a more plausible role when
one or more of those assumptions fails:
\begin{itemize}
    \item non-smooth sources distributed throughout the domain,
    \item close-to-source quantities where low-order far-field moments are
    insufficient,
    \item many interacting parameters rather than one or two frequency-like
    parameters,
    \item tolerance targets closer to a few percent than to machine precision,
    \item nonlinear wave equations without a Green's function or
    superposition.
\end{itemize}

This list should guide the interpretation of the thesis. The linear FD/PML
study is not the final application where learned transfer is most competitive.
It is the baseline difficulty test that a learned method must pass before the
nonlinear or high-parameter setting is credible.

% ============================================================
\section{Training cost and amortisation}
\label{sec:discussion-amortisation}
% ============================================================

The cost of training is a first-order limitation. The current two-dimensional
datasets use $9600$ exact FD/PML solves per frequency pair. Such a cost can be
justified only if the trained model is reused across many subsequent solves or
if the trained model enables a regime that classical methods cannot handle.

\subsection{The fixed-operator requirement}
\label{subsec:discussion-fixed-operator}

The most favourable linear use case is a fixed operator with many right-hand
sides. The medium, grid, boundary treatment, and frequency pair must remain
unchanged enough that the trained map can be reused. If the medium changes
after each solve, as in many full-waveform inversion workflows, the amortised
value of a fixed trained model becomes much weaker.

This does not make the training programme invalid, but it narrows the
practical claim. The current method is most plausible for repeated fixed-
medium query settings, design libraries, monitoring problems with a stable
background operator, or offline surrogate construction in which many future
right-hand sides are expected.

\subsection{Data efficiency is part of the solver claim}
\label{subsec:discussion-data-efficiency}

Solver acceleration cannot be judged only by per-solve iteration count. The
total cost includes data generation, training, evaluation, and any classical
preconditioner still used during deployment. A warm start that saves one or
two FGMRES iterations is not economically meaningful if it required thousands
of training solves and is reused only a small number of times.

The thesis results therefore suggest a practical reporting standard for future
work: every learned solver component should report not only residual reduction
and iteration count, but also training solves, high-fidelity labels required,
and the break-even number of deployment solves. This would make the
amortisation assumption visible rather than implicit.

\subsection{Source priors may reduce the cost}
\label{subsec:discussion-source-priors}

The present datasets draw sources from a broad interior distribution. Real
applications often have much sharper source priors, such as antenna edges,
localized emitters, repeated sensor layouts, or clustered forcing regions.
Restricting the source distribution can reduce the complexity of the learned
map and therefore the number of training solves required. This is not a result
demonstrated in the thesis, but it is a concrete way to test whether the cost
barrier can be softened without changing the solver method.

\subsection{The nonlinear amortisation argument}
\label{subsec:discussion-nonlinear-amortisation}

The most compelling long-term cost argument is not the linear homogeneous
case. In nonlinear wave problems there is no Green's function, no
superposition, and no direct quadrature over precomputed linear responses.
Newton-type methods may also require a good initial state to converge at all,
not merely to converge faster. In that regime, a learned warm start can become
a feasibility mechanism rather than a marginal iteration-saving device.

The linear study is therefore defensible as a baseline of difficulty. If
learned frequency transfer cannot produce useful structure in the linear case,
then a nonlinear extension has little chance. The open scientific question is
whether the classical-versus-learned difficulty gap shrinks as nonlinearity
grows. Classical Green's-function methods become much less applicable; neural
surrogates may become only moderately harder because nonlinearity is already
part of the model class.

% ============================================================
\section{Limitations and threats to validity}
\label{sec:discussion-limitations}
% ============================================================

\subsection{Linear homogeneous physics}
\label{subsec:discussion-limit-linear}

The main experiments use homogeneous linear Helmholtz systems. This makes the
operator-theoretic analysis clean and the exact-solve datasets reproducible,
but it also gives classical methods the maximum possible advantage. The thesis
does not demonstrate that learned frequency transfer is competitive with
Green's-function or low-rank methods in this regime. It demonstrates a
residual-amplification mechanism and a warm-start effect inside a controlled
linear baseline.

\subsection{Warm start rather than preconditioner}
\label{subsec:discussion-limit-warmstart}

The two-dimensional learned model sets $x_0$ and then leaves CSL-FGMRES to
solve the system. It is not applied as $M^{-1}$ inside each Krylov iteration.
The results therefore establish warm-start behaviour only. A learned
preconditioner claim would require residual-to-correction training on Krylov
states and evaluation inside FGMRES.

\subsection{Limited sample size and frequency coverage}
\label{subsec:discussion-limit-samples}

The main two-dimensional beta-$0.3$ solver evaluation uses $10$ test
right-hand sides per pair and three dyadic pairs:
$16\to32$, $32\to64$, and $64\to128$. These runs are sufficient to support the
reported qualitative ordering, but they do not provide tight confidence
intervals over source distributions, random training seeds, or alternative
frequency ratios.

\subsection{Single main CSL damping}
\label{subsec:discussion-limit-beta}

The final discussion focuses on $\beta=0.3$. Earlier beta-$0.1$ runs show
compatible qualitative trends, but the thesis does not contain a systematic
modern beta sweep for the flux-full model. Any claim about the method should
therefore be read at fixed CSL damping. The learned warm start and the
classical preconditioner interact, and changing $\beta$ can change the
residual modes that matter.

\subsection{Non-normal PML operator}
\label{subsec:discussion-limit-nonnormal}

The Dirichlet modal analysis relies on an orthonormal eigenbasis. The PML
operator is complex-valued and non-normal, so eigenvalues alone do not control
transient behaviour or residual amplification. Pseudospectral effects can
make a residual component large even when eigenvalue plots look benign
\cite{TrefethenEmbree2005}. This is why the 1D PML and 2D FD/PML sections rely on
direct residual diagnostics rather than claiming a full modal theory.

\subsection{Architecture coverage}
\label{subsec:discussion-limit-architecture}

The two-dimensional solver-facing model is a TransferUNet. The older
architecture comparison shows that the dilated CNN and U-Net are broadly
comparable in field-error scaling, but the thesis does not provide a
two-dimensional solver-facing comparison against Fourier neural operators,
transformers, or graph-based neural operators. The residual-amplification
diagnostic is architecture-independent, but the achieved residual distribution
is not proven architecture-independent.

\subsection{Tool-first research framing}
\label{subsec:discussion-limit-tool-first}

The project began as an investigation of whether convolutional neural networks
could learn frequency transfer. That is a tool-first research mode. The risk
is that the tool may be applied to a problem where classical structure is
already superior. The scientific value of the thesis is that the tool-first
investigation exposed a problem-first mechanism: residual amplification is the
constraint that any learned or classical transfer approximation must satisfy
before it becomes solver-useful.

% ============================================================
\section{Future research directions}
\label{sec:discussion-future}
% ============================================================

\subsection{Residual-aware warm-start training}
\label{subsec:discussion-future-residual}

The most direct next experiment is to replace pure field loss with a
residual-aware objective:
\begin{equation}
    \mathcal{L}_{\mathrm{res}}
    =
    \frac{\|A_H \hat{u}_H - b\|_2^2}{\|b\|_2^2+\varepsilon}.
    \label{eq:discussion-residual-loss}
\end{equation}
This loss asks the network to satisfy the discrete high-frequency equation,
not only to match the solution field. The current results suggest that this
objective should be evaluated by initial true residual, initial
preconditioned residual, final true residual, and residual area under the
FGMRES curve, not by field error alone.

Generic residual loss may still be insufficient, as the 1D learned correction
experiments indicate. The stronger version is a weighted residual loss that
emphasises modes surviving late in CSL-FGMRES:
\begin{equation}
    \mathcal{L}_{\mathrm{late}}
    =
    \sum_k w_k
    \left|
    \left\langle v_k, A_H e_\theta \right\rangle
    \right|^2,
    \label{eq:discussion-late-mode-loss}
\end{equation}
where $w_k$ is large for residual modes that remain after several CSL-FGMRES
iterations. In the Dirichlet setting the weights can be defined spectrally. In
the PML setting they should be defined through measured Krylov residuals or
through a pseudospectral proxy.

\subsection{Exact-coarse imitation}
\label{subsec:discussion-future-exact-coarse}

The exact low solve plus $T_\uparrow$ residual gate is the current proof of
mechanism. The most targeted learned extension is therefore not simply
``train on residuals'' but imitate the accepted modal action of this exact
coarse correction. The target should match scale, phase, and the subset of
modes that reduce the residual after CSL has done its part.

A clean 1D ablation would compare five objectives:
\begin{itemize}
    \item field loss,
    \item unweighted residual loss,
    \item spectral residual loss,
    \item late-CSL weighted residual loss,
    \item exact-coarse imitation loss.
\end{itemize}
The win condition should be fewer CSL-FGMRES iterations or lower final true
residual at fixed iteration budget. A prettier field is not the win condition.

\subsection{Residual-to-correction data from Krylov states}
\label{subsec:discussion-future-rescorr}

The two-dimensional preconditioner extension should train on the object seen
inside Krylov: residuals and corrections, not complete solution fields. The
natural dataset consists of tuples
\[
    (r_k, e_k),
    \qquad
    A_H e_k = r_k,
\]
or practical variants based on CSL-corrected residuals. These tuples should be
generated from actual CSL-FGMRES states under the full FD/PML operator, with
the PML region included. Training on such data would align the deployment
distribution with the training distribution.

The learned correction should be used conservatively. A residual gate can
accept a proposed correction only if it decreases the true residual, the
preconditioned residual, or both. Such a gate is computationally more
expensive than a blind neural step, but it preserves the main lesson of the
thesis: learned components must be judged by the residual the solver sees.

\subsection{Hybrid physics-enhanced surrogates}
\label{subsec:discussion-future-peds}

The most promising long-term architecture is hybrid rather than purely
neural. A classical low-fidelity or coarse solver can handle the part of the
problem already understood, while a neural network learns the correction that
the classical model misses. Physics-enhanced deep surrogates follow this
template by embedding a low-fidelity physics simulator inside an end-to-end
trained surrogate and have been reported to reduce data requirements
substantially for PDE surrogate tasks \cite{PestourieEtAl2023PEDS}.

For frequency transfer, the low-fidelity component could be a coarse-grid
solve, a quadrature Green's-function approximation, a low-rank surrogate, or a
CSL correction. The network would then learn the residual correction rather
than the whole solution map. This design answers the training-cost criticism
more directly than a standalone neural field predictor, because it preserves
as much classical structure as possible.

\subsection{Nonlinear wave extensions}
\label{subsec:discussion-future-nonlinear}

The nonlinear extension is the most important application-level direction.
For nonlinear frequency-domain wave equations, superposition fails and Green's
function convolution is unavailable. Classical perturbative methods can handle
weak nonlinearities, but strongly nonlinear regimes may require good initial
states for Newton or continuation methods to converge. A learned frequency
transfer map could then supply a feasibility-enabling initial guess rather
than a marginal acceleration.

The thesis results define the entry test for that programme. The learned model
must first pass the linear baseline, including residual compatibility. The
next question is whether adding nonlinearity makes classical methods much
harder while making the learned map only moderately harder. If that gap
shrinks as nonlinearity grows, learned frequency transfer becomes a more
natural tool.

\subsection{Source-conditioned and prior-aware models}
\label{subsec:discussion-future-source}

The current two-dimensional flux-full model maps low-frequency fields to
high-frequency fields without explicit source channels in the main trained
model. The datasets store source information, and a source-conditioned model
is a natural extension. Source channels can help distinguish fields that look
similar at low frequency but require different high-frequency corrections near
the source. Source priors can also reduce training cost by focusing the model
on physically relevant forcing distributions.

\subsection{Architecture-independence tests}
\label{subsec:discussion-future-architecture}

The residual-amplification threshold predicts that different architectures
with matched residual spectra should behave similarly as warm starts, even if
their field errors differ. A two-dimensional comparison against a Fourier
neural operator would test whether the current observations are U-Net-specific
or operator-level. The experiment should match training data, loss, source
distribution, and solver evaluation protocol. The metric should be residual
behaviour, not only field RelL2 \cite{LiEtAl2021FNO}.

\subsection{Three-dimensional deployment}
\label{subsec:discussion-future-3d}

Three-dimensional Helmholtz problems are the real target in seismic,
photonic, and acoustic applications. The computational savings from a useful
warm start or correction could be much larger there, but the residual-safety
threshold is also expected to be sharper because the systems are larger and
more indefinite. A three-dimensional extension is therefore not a routine
scaling exercise. It is a new test of whether the learned correction can meet
a stricter operator-amplification constraint.

% ============================================================
\section{Closing interpretation}
\label{sec:discussion-closing}
% ============================================================

The thesis began with the observation that adjacent Helmholtz frequencies are
not independent. A low-frequency solution contains information about the
corresponding high-frequency solution, and a neural network can learn part of
that map. The results show that this observation is true but incomplete. A
Krylov solver does not consume visual similarity, field RelL2, or
cross-frequency coherence directly. It consumes residuals.

The core contribution is therefore a change in the question. The important
question is not whether a network can predict the high-frequency field. The
important question is whether the error in that prediction lies in directions
that the discrete high-frequency operator and the CSL-preconditioned Krylov
process can tolerate. The answer is sometimes yes, often only after
residual-aware selection, and not yet strongly enough to claim a learned
preconditioner.

This is a useful endpoint for a thesis because it turns an initially broad
neural-solver idea into a falsifiable numerical-analysis statement. Learned
frequency transfer is solver-useful when the induced residual clears the
operator-amplification barrier. Field loss alone does not enforce that
condition. The next method should therefore be residual-aware, PML-compatible,
and hybrid with the classical solvers that already solve the parts of the
linear problem that learning should not relearn.

\todo[inline]{Bibliography checks before final compile: add or verify entries
for \texttt{Erlangga2004}, \texttt{Engquist2011},
\texttt{ParksEtAl2006}, \texttt{PestourieEtAl2023PEDS}, and
\texttt{LiEtAl2021FNO}. The PEDS article is Pestourie, Mroueh, Rackauckas,
Das, and Johnson, ``Physics-enhanced deep surrogates for partial differential
equations,'' Nature Machine Intelligence 5, 1458 to 1465, 2023,
DOI 10.1038/s42256-023-00761-y.}

\todo[inline]{Number checks before submission: verify the 2D beta-$0.3$ table
against the final CSVs under the frozen checkpoint evaluation. The values used
here match the current \texttt{chapter\_results\_after\_kees.tex} draft and
the local 2026-05-10/2026-05-18 notes.}

